\documentclass{amsart}
% For dealing with references we use the comment environment
\usepackage{verbatim}
\newenvironment{reference}{\comment}{\endcomment}
%\newenvironment{reference}{}{}
\newenvironment{slogan}{\comment}{\endcomment}
\newenvironment{history}{\comment}{\endcomment}

% For commutative diagrams we use Xy-pic
\usepackage[all]{xy}

% We use 2cell for 2-commutative diagrams.
\xyoption{2cell}
\UseAllTwocells

% We use multicol for the list of chapters between chapters
\usepackage{multicol}

% This is generally recommended for better output
\usepackage{lmodern}
\usepackage[T1]{fontenc}

% For cross-file-references

% Package for hypertext links:
\usepackage{hyperref}

% For any local file, say "hello.tex" you want to link to please

% Theorem environments.
%
\theoremstyle{plain}
\newtheorem{theorem}[subsection]{Theorem}
\newtheorem{proposition}[subsection]{Proposition}
\newtheorem{lemma}[subsection]{Lemma}

\theoremstyle{definition}
\newtheorem{definition}[subsection]{Definition}
\newtheorem{example}[subsection]{Example}
\newtheorem{exercise}[subsection]{Exercise}
\newtheorem{situation}[subsection]{Situation}

\theoremstyle{remark}
\newtheorem{remark}[subsection]{Remark}
\newtheorem{remarks}[subsection]{Remarks}

\numberwithin{equation}{subsection}

% Macros
%
\def\lim{\mathop{\mathrm{lim}}\nolimits}
\def\colim{\mathop{\mathrm{colim}}\nolimits}
\def\Spec{\mathop{\mathrm{Spec}}}
\def\Hom{\mathop{\mathrm{Hom}}\nolimits}
\def\Ext{\mathop{\mathrm{Ext}}\nolimits}
\def\SheafHom{\mathop{\mathcal{H}\!\mathit{om}}\nolimits}
\def\SheafExt{\mathop{\mathcal{E}\!\mathit{xt}}\nolimits}
\def\Sch{\mathit{Sch}}
\def\Mor{\mathop{\mathrm{Mor}}\nolimits}
\def\Ob{\mathop{\mathrm{Ob}}\nolimits}
\def\Sh{\mathop{\mathit{Sh}}\nolimits}
\def\NL{\mathop{N\!L}\nolimits}
\def\CH{\mathop{\mathrm{CH}}\nolimits}
\def\proetale{{pro\text{-}\acute{e}tale}}
\def\etale{{\acute{e}tale}}
\def\QCoh{\mathit{QCoh}}
\def\Ker{\mathop{\mathrm{Ker}}}
\def\Im{\mathop{\mathrm{Im}}}
\def\Coker{\mathop{\mathrm{Coker}}}
\def\Coim{\mathop{\mathrm{Coim}}}

% Boxtimes
%
\DeclareMathSymbol{\boxtimes}{\mathbin}{AMSa}{"02}

%
% Macros for moduli stacks/spaces
%
\def\QCohstack{\mathcal{QC}\!\mathit{oh}}
\def\Cohstack{\mathcal{C}\!\mathit{oh}}
\def\Spacesstack{\mathcal{S}\!\mathit{paces}}
\def\Quotfunctor{\mathrm{Quot}}
\def\Hilbfunctor{\mathrm{Hilb}}
\def\Curvesstack{\mathcal{C}\!\mathit{urves}}
\def\Polarizedstack{\mathcal{P}\!\mathit{olarized}}
\def\Complexesstack{\mathcal{C}\!\mathit{omplexes}}
% \Pic is the operator that assigns to X its picard group, usage \Pic(X)
% \Picardstack_{X/B} denotes the Picard stack of X over B
% \Picardfunctor_{X/B} denotes the Picard functor of X over B
\def\Pic{\mathop{\mathrm{Pic}}\nolimits}
\def\Picardstack{\mathcal{P}\!\mathit{ic}}
\def\Picardfunctor{\mathrm{Pic}}
\def\Deformationcategory{\mathcal{D}\!\mathit{ef}}

\setlength{\emergencystretch}{3em}
\begin{document}
\title[Stacks Project, Tag 095U]{287 --- Relative cohomological-dimension bound twice the fibre dimension for proper morphisms}
\maketitle
\noindent Copyright (C) 2005--2025 Johan de Jong. Stacks Project authors. Source: \href{https://stacks.math.columbia.edu/tag/095U}{Tag 095U}.

\noindent Editorial difficulty note (not author text). Difficulty H2: The proof must control arbitrary torsion sheaves by dimension induction, including singular and nonnormal schemes. Normalization produces a lower-dimensional cokernel, while the graph of a rational map to the projective line produces a derived error supported in codimension at least two; Leray then gives the sharp twice-dimension bound..

\noindent Editorial provenance note (not author text). History: excerpt from revision \texttt{a0444\allowbreak 6e57e\allowbreak c1fbc\allowbreak 252a8\allowbreak 71afc\allowbreak ec775\allowbreak 2fb28\allowbreak 07b14}, etale-cohomology.tex, lines 18785--18880. Reference presentation was adapted; original-author context and core support blocks were retained or added. The mathematical wording is unchanged. Licensed under GNU FDL 1.2 or later; no invariant sections or cover texts. See ../licenses/COPYING. Context blocks reproduce author definitions and supporting results; they are not additional counted problems.

\begin{lemma}
\label{lemma-cohomological-dimension-proper}
Let $f : X \to Y$ be a proper morphism of schemes
all of whose fibres have dimension $\leq n$.
Then for any abelian torsion sheaf $\mathcal{F}$ on $X_\etale$
we have $R^qf_*\mathcal{F} = 0$ for $q > 2n$.
\end{lemma}

\begin{proof}
We will prove this by induction on $n$ for all proper morphisms.

\medskip\noindent
If $n = 0$, then $f$ is a finite morphism
(More on Morphisms, Lemma \href{https://stacks.math.columbia.edu/tag/02LS}{lemma characterize finite (Tag 02LS)})
and the
result is true by Proposition \href{https://stacks.math.columbia.edu/tag/03QP}{proposition finite higher direct image zero (Tag 03QP)}.

\medskip\noindent
If $n > 0$, then using Lemma \href{https://stacks.math.columbia.edu/tag/0DDF}{lemma proper base change stalk (Tag 0DDF)}
we see that it suffices to prove $H^i_\etale(X, \mathcal{F}) = 0$
for $i > 2n$ and $X$ a proper scheme, $\dim(X) \leq n$
over an algebraically closed field $k$
and $\mathcal{F}$ is a torsion abelian sheaf on $X$.

\medskip\noindent
If $n = 1$ this follows from Theorem \href{https://stacks.math.columbia.edu/tag/03RT}{theorem vanishing curves (Tag 03RT)}.
Assume $n > 1$. By Proposition \href{https://stacks.math.columbia.edu/tag/03SI}{proposition topological invariance (Tag 03SI)}
we may replace $X$ by its reduction.
Let $\nu : X^\nu \to X$ be the normalization.
This is a surjective birational finite morphism
(see Varieties, Lemma \href{https://stacks.math.columbia.edu/tag/0BXR}{lemma normalization locally algebraic (Tag 0BXR)})
and hence an isomorphism over a dense open $U \subset X$
(Morphisms, Lemma \href{https://stacks.math.columbia.edu/tag/0BAC}{lemma birational birational (Tag 0BAC)}).
Then we see that
$c : \mathcal{F} \to \nu_*\nu^{-1}\mathcal{F}$ is injective
(as $\nu$ is surjective) and an isomorphism over $U$.
Denote $i : Z \to X$ the inclusion of the complement of $U$.
Since $U$ is dense in $X$ we have $\dim(Z) < \dim(X) = n$.
By Proposition \href{https://stacks.math.columbia.edu/tag/04CA}{proposition closed immersion pushforward (Tag 04CA)} have
$\Coker(c) = i_*\mathcal{G}$ for some
abelian torsion sheaf $\mathcal{G}$ on $Z_\etale$.
Then $H^q_\etale(X, \Coker(c)) = H^q_\etale(Z, \mathcal{G})$
(by Proposition \href{https://stacks.math.columbia.edu/tag/03QP}{proposition finite higher direct image zero (Tag 03QP)}
and the Leray spectral sequence) and by induction hypothesis we conclude that
the cokernel of $c$ has cohomology in degrees $\leq 2(n - 1)$.
Thus it suffices to prove the result for $\nu_*\nu^{-1}\mathcal{F}$.
As $\nu$ is finite this reduces us to showing
that $H^i_\etale(X^\nu, \nu^{-1}\mathcal{F})$ is
zero for $i > 2n$. This case is treated in the next paragraph.

\medskip\noindent
Assume $X$ is integral normal proper scheme over $k$ of dimension $n$.
Choose a nonconstant rational function $f$ on $X$. The graph
$X' \subset X \times \mathbf{P}^1_k$ of $f$ sits into a diagram
$$
X \xleftarrow{b} X' \xrightarrow{f} \mathbf{P}^1_k
$$
Observe that $b$ is an isomorphism over an open subscheme
$U \subset X$ whose complement is a closed subscheme
$Z \subset X$ of codimension $\geq 2$. Namely, $U$ is the
domain of definition of $f$ which contains all codimension $1$
points of $X$, see
Morphisms, Lemmas \href{https://stacks.math.columbia.edu/tag/0A1Y}{lemma rational map from reduced to separated (Tag 0A1Y)}
and \href{https://stacks.math.columbia.edu/tag/0BX7}{lemma extend across (Tag 0BX7)}
(combined with Serre's criterion for normality, see
Properties, Lemma \href{https://stacks.math.columbia.edu/tag/0345}{lemma criterion normal (Tag 0345)}).
Moreover the fibres of $b$ have dimension $\leq 1$ (as closed subschemes
of $\mathbf{P}^1$). Hence $R^ib_*b^{-1}\mathcal{F}$ is nonzero only
if $i \in \{0, 1, 2\}$ by induction. Choose a distinguished triangle
$$
\mathcal{F} \to Rb_*b^{-1}\mathcal{F} \to Q \to \mathcal{F}[1]
$$
Using that $\mathcal{F} \to b_*b^{-1}\mathcal{F}$ is injective
as before and using what we just said, we see that $Q$ has nonzero
cohomology sheaves only in degrees $0, 1, 2$ sitting on $Z$.
Moreover, these cohomology sheaves are torsion by
Lemma \href{https://stacks.math.columbia.edu/tag/0DDD}{lemma torsion direct image (Tag 0DDD)}.
By induction and the spectral sequence of
Derived Categories, Lemma \href{https://stacks.math.columbia.edu/tag/015J}{lemma two ss complex functor (Tag 015J)}
we see that $H^i(X, Q)$ is zero for
$i > 2 + 2\dim(Z) \leq 2 + 2(n - 2) = 2n - 2$. Thus it suffices
to prove that $H^i(X', b^{-1}\mathcal{F}) = 0$ for
$i > 2n$. At this point we use the morphism
$$
f : X' \to \mathbf{P}^1_k
$$
whose fibres have dimension $< n$. Hence by induction we see that
$R^if_*b^{-1}\mathcal{F} = 0$ for $i > 2(n - 1)$.
We conclude by the Leray spectral sequence
$$
H^i(\mathbf{P}^1_k, R^jf_*b^{-1}\mathcal{F})
\Rightarrow
H^{i + j}(X', b^{-1}\mathcal{F})
$$
and the fact that $\dim(\mathbf{P}^1_k) = 1$.
\end{proof}
\end{document}
